\section{NFV-based benchmarking alpha as a composite object}
\label{sec:appendix_bv_alpha_interpretation}

This appendix section provides an auxiliary interpretation of NFV-based pricing errors for ex-post benchmarking.
Unlike the baseline structural restriction in Equation~\eqref{eq:deal_pricing}, the decomposition below allows for a nonzero conditional abnormal component at deal inception and is therefore used only as an interpretive device, not to define the structural estimation target.

\begin{remark}[Measured alpha versus conditional abnormal performance]
	\label{rem:alpha_decomp}
	Fix a deal $j$ of fund $i$ and suppose $\tau\le t_{i,j}^{\mathrm{Inv}}<t_{i,j}^{\mathrm{Future}}$.
	Define the deal-level conditional abnormal component at inception by
	\[
	a_{i,j}^{\mathrm{skill}}
	:=
	\mathbb{E}\!\left[
	\delta_{i,j}\mid
	\mathcal{F}_{t_{i,j}^{\mathrm{Inv}}}
	\right].
	\]
	Write
	\[
	\delta_{i,j}
	=
	a_{i,j}^{\mathrm{skill}}
	+
	u_{i,j},
	\qquad
	\mathbb{E}\!\left[
	u_{i,j}\mid
	\mathcal{F}_{t_{i,j}^{\mathrm{Inv}}}
	\right]
	=
	0.
	\]
	Then the conditional expectation of the NFV-type moment can be decomposed as
	\begin{align}
		\label{eq:alpha_decomp_general}
		\mathbb{E}\!\left[
		\frac{\Psi_{t_{i,j}^{\mathrm{Future}}}}{\Psi_\tau}\,
		\delta_{i,j}
		\,\middle|\,
		\mathcal{F}_\tau
		\right]
		&=
		\mathbb{E}\!\left[
		\frac{\Psi_{t_{i,j}^{\mathrm{Inv}}}}{\Psi_\tau}
		\left(
		\mathbb{E}\!\left[
		\frac{\Psi_{t_{i,j}^{\mathrm{Future}}}}{\Psi_{t_{i,j}^{\mathrm{Inv}}}}
		\,\middle|\,
		\mathcal{F}_{t_{i,j}^{\mathrm{Inv}}}
		\right]
		a_{i,j}^{\mathrm{skill}}
		\right)
		\,\middle|\,
		\mathcal{F}_\tau
		\right]
		\nonumber\\
		&\qquad
		+
		\mathbb{E}\!\left[
		\frac{\Psi_{t_{i,j}^{\mathrm{Inv}}}}{\Psi_\tau}\,
		\mathrm{Cov}\!\left(
		\frac{\Psi_{t_{i,j}^{\mathrm{Future}}}}{\Psi_{t_{i,j}^{\mathrm{Inv}}}},
		\delta_{i,j}
		\,\middle|\,
		\mathcal{F}_{t_{i,j}^{\mathrm{Inv}}}
		\right)
		\,\middle|\,
		\mathcal{F}_\tau
		\right].
	\end{align}
	
	If, in addition, the gross risk-free return over the holding period is defined by
	\[
	\mathbb{E}\!\left[
	\frac{\Psi_{t_{i,j}^{\mathrm{Future}}}}{\Psi_{t_{i,j}^{\mathrm{Inv}}}}
	\,\middle|\,
	\mathcal{F}_{t_{i,j}^{\mathrm{Inv}}}
	\right]
	=
	\frac{1}{R_{f,t_{i,j}^{\mathrm{Inv}}\to t_{i,j}^{\mathrm{Future}}}},
	\]
	then Equation~\eqref{eq:alpha_decomp_general} becomes
	\begin{align}
		\label{eq:alpha_decomp}
		\underbrace{
			\mathbb{E}\!\left[
			\frac{\Psi_{t_{i,j}^{\mathrm{Future}}}}{\Psi_\tau}\,
			\delta_{i,j}
			\,\middle|\,
			\mathcal{F}_\tau
			\right]
		}_{\text{Measured NFV alpha}}
		&=
		\underbrace{
			\mathbb{E}\!\left[
			\frac{\Psi_{t_{i,j}^{\mathrm{Inv}}}}{\Psi_\tau}\,
			\frac{a_{i,j}^{\mathrm{skill}}}{R_{f,t_{i,j}^{\mathrm{Inv}}\to t_{i,j}^{\mathrm{Future}}}}
			\,\middle|\,
			\mathcal{F}_\tau
			\right]
		}_{\text{Conditional abnormal component}}
		\nonumber\\
		&\qquad
		+
		\underbrace{
			\mathbb{E}\!\left[
			\frac{\Psi_{t_{i,j}^{\mathrm{Inv}}}}{\Psi_\tau}\,
			\mathrm{Cov}\!\left(
			\frac{\Psi_{t_{i,j}^{\mathrm{Future}}}}{\Psi_{t_{i,j}^{\mathrm{Inv}}}},
			\delta_{i,j}
			\,\middle|\,
			\mathcal{F}_{t_{i,j}^{\mathrm{Inv}}}
			\right)
			\,\middle|\,
			\mathcal{F}_\tau
			\right]
		}_{\text{Timing Risk Premium}}.
	\end{align}
	Thus, an NFV-based ex-post alpha generally mixes a conditional abnormal component with a timing-related covariance component.
	Even when $a_{i,j}^{\mathrm{skill}}=0$, the measured NFV alpha can still be nonzero because of the Timing Risk Premium.
\end{remark}